\documentclass[11pt]{article}
\usepackage[margin=1.02in]{geometry}
\usepackage{amsmath,amssymb,amsthm,mathtools}
\usepackage{graphicx}
\usepackage{xcolor}
\usepackage[hidelinks]{hyperref}
\usepackage{microtype}

\newtheorem{theorem}{Theorem}
\newtheorem{remark}{Remark}
\newcommand{\C}{\mathbb C}
\newcommand{\Tr}{\operatorname{Tr}}
\newcommand{\diag}{\operatorname{diag}}
\newcommand{\Cp}{\mathcal C_p}
\newcommand{\Sp}{\mathsf S_p}

\definecolor{statusblue}{RGB}{31,78,121}

\title{Norm Compression for Scaled-Unitary Blocks}
\author{fa\_banach\_001 / agent\_lane\_00}
\date{19 July 2026}

\begin{document}
\maketitle

\noindent
\fcolorbox{statusblue}{statusblue!6}{%
  \begin{minipage}{0.92\textwidth}
  \textbf{Status: candidate partial result, likely valid.}
  Audenaert's general norm-compression conjecture remains open in this packet.
  The theorem below proves all finite Schatten exponents for arbitrary
  collections of common-size blocks that are scalar multiples of unitaries.
  \end{minipage}}

\section{Source conjecture}

Koenraad M. R. Audenaert's paper \emph{On a Norm Compression Inequality for
$2\times N$ Partitioned Block Matrices}, arXiv:math/0702186, formulates
Conjecture~1 on PDF page~2 \cite{Audenaert}.  For
\[
 T=\begin{pmatrix}A_1&A_2&\cdots&A_N\\ B_1&B_2&\cdots&B_N\end{pmatrix},
 \qquad
 \Cp(T)=\begin{pmatrix}
 \lVert A_1\rVert_p&\cdots&\lVert A_N\rVert_p\\
 \lVert B_1\rVert_p&\cdots&\lVert B_N\rVert_p
 \end{pmatrix},
\]
it asks whether
\begin{align}
 \lVert T\rVert_p&\geq \lVert\Cp(T)\rVert_p,
    &&1\leq p\leq2, \label{eq:lower}\\
 \lVert T\rVert_p&\leq \lVert\Cp(T)\rVert_p,
    &&p\geq2. \label{eq:upper}
\end{align}
Here every norm is the Schatten norm at the displayed exponent.

\begin{figure}[ht]
  \centering
  \includegraphics[width=\textwidth]{figures/open_problem_crop.png}
  \caption{The definition and Conjecture~1 as rendered from page~2 of the
  source paper.}
\end{figure}

\clearpage
\section{Definitions and scope}

For a finite matrix $X$, its Schatten $p$-norm is
\[
 \lVert X\rVert_p=\bigl(\Tr |X|^p\bigr)^{1/p},\qquad 1\leq p<\infty.
\]
We treat the subcase in which all blocks are $d\times d$ and
\[
 A_i=a_iU_i,\qquad B_i=b_iV_i,
 \qquad a_i,b_i\geq0,\quad U_i,V_i\in\mathrm U(d).
\]
The unitaries may depend arbitrarily on $i$.  Thus the blocks may be full
rank, non-proportional, and noncommuting.  This differs from the rank-one and
proportional-block families handled in the source.

\section{Proof intuition}

Scaled unitaries make the two diagonal blocks of $TT^*$ scalar.  All remaining
geometry is contained in one cross term
\[
 Z=\sum_i a_ib_iU_iV_i^*.
\]
After taking a singular-value decomposition of $Z$, the $2d\times2d$ matrix
$TT^*$ splits into $d$ scalar $2\times2$ problems.  If $s_j$ is a singular
value of $Z$, the triangle inequality gives
$s_j\leq c:=\sum_i a_ib_i$.  The scalar compression is exactly the perfectly
aligned endpoint $s=c$, repeated $d$ times.

The sum of the $q$th powers of the two eigenvalues of
$\bigl(\begin{smallmatrix}\alpha&s\\s&\beta\end{smallmatrix}\bigr)$ is
decreasing in $s$ for $q\leq1$ and increasing for $q\geq1$.  Substituting
$q=p/2$ produces the desired reversal at $p=2$.  The unrestricted conjecture
remains open because general blocks do not make the diagonal blocks of $TT^*$
scalar, so this singular-value decoupling is unavailable.

\section{The partial theorem}

\begin{theorem}[Scaled-unitary norm compression]
\label{thm:main}
Let $N,d\geq1$, let $a_i,b_i\geq0$, and let $U_i,V_i\in\mathrm U(d)$ for
$1\leq i\leq N$.  Form
\[
 T=\begin{pmatrix}
 a_1U_1&\cdots&a_NU_N\\
 b_1V_1&\cdots&b_NV_N
 \end{pmatrix}.
\]
Then Audenaert's norm-compression inequalities \eqref{eq:lower} and
\eqref{eq:upper} hold for every $1\leq p<\infty$.
\end{theorem}

\begin{proof}
Fix $p$ and write $q=p/2$.  Set
\[
 \alpha=\sum_{i=1}^N a_i^2,\qquad
 \beta=\sum_{i=1}^N b_i^2,\qquad
 c=\sum_{i=1}^N a_ib_i,
 \qquad
 Z=\sum_{i=1}^N a_ib_iU_iV_i^*.
\]
The unitary hypotheses give
\begin{equation}
 TT^*=\begin{pmatrix}
 \alpha I_d&Z\\ Z^*&\beta I_d
 \end{pmatrix}.                                      \label{eq:gram}
\end{equation}
Let $s_1,\ldots,s_d$ be the singular values of $Z$.  If
$Z=L\diag(s_1,\ldots,s_d)R^*$ is a singular-value decomposition, conjugating
\eqref{eq:gram} by $\diag(L,R)$ and permuting coordinates shows that $TT^*$
is unitarily equivalent to the direct sum
\[
 \bigoplus_{j=1}^d
 \begin{pmatrix}\alpha&s_j\\s_j&\beta\end{pmatrix}.
\]
Moreover,
\begin{equation}
 0\leq s_j\leq\lVert Z\rVert_\infty
 \leq\sum_{i=1}^N a_ib_i\lVert U_iV_i^*\rVert_\infty=c.       \label{eq:sbound}
\end{equation}

For $0\leq s\leq c$, define
\[
 \lambda_\pm(s)=\frac{\alpha+\beta
 \pm\sqrt{(\alpha-\beta)^2+4s^2}}{2},
 \qquad
 F_q(s)=\lambda_+(s)^q+\lambda_-(s)^q.
\]
Cauchy--Schwarz gives $c^2\leq\alpha\beta$, so these are the two
nonnegative eigenvalues throughout the required interval.  The preceding
unitary equivalence yields
\begin{equation}
 \lVert T\rVert_p^p=\Tr((TT^*)^q)=\sum_{j=1}^dF_q(s_j).       \label{eq:full}
\end{equation}

For $s>0$ at points where both eigenvalues are positive, direct
differentiation gives
\begin{equation}
 F_q'(s)=
 \frac{2qs}{\sqrt{(\alpha-\beta)^2+4s^2}}
 \left(\lambda_+(s)^{q-1}-\lambda_-(s)^{q-1}\right).        \label{eq:derivative}
\end{equation}
Thus $F_q$ is nonincreasing when $0<q\leq1$ and nondecreasing when $q\geq1$.
The same conclusion at $s=0$ and when $\lambda_-(s)=0$ follows by continuity.
At $q=1$, $F_1(s)=\alpha+\beta$ is constant.

It remains to identify the endpoint $s=c$ with the norm compression.  Since
the singular values of each unitary equal one,
\[
 \lVert a_iU_i\rVert_p=d^{1/p}a_i,
 \qquad
 \lVert b_iV_i\rVert_p=d^{1/p}b_i.
\]
Consequently
\[
 \Cp(T)\Cp(T)^*=d^{2/p}
 \begin{pmatrix}\alpha&c\\c&\beta\end{pmatrix},
\]
and, because $2q/p=1$,
\begin{equation}
 \lVert\Cp(T)\rVert_p^p
 =\Tr\bigl((\Cp(T)\Cp(T)^*)^q\bigr)
 =dF_q(c).                                                \label{eq:compressed}
\end{equation}

If $1\leq p\leq2$, then $q\leq1$; equations
\eqref{eq:sbound}--\eqref{eq:compressed} give
\[
 \lVert T\rVert_p^p=\sum_jF_q(s_j)\geq dF_q(c)
 =\lVert\Cp(T)\rVert_p^p.
\]
If $p\geq2$, the monotonicity reverses and gives the opposite inequality.
Taking $p$th roots proves the theorem.
\end{proof}

\begin{remark}
At $p=2$ there is equality for every choice of blocks.  Away from $p=2$, the
proof also shows equality whenever every singular value of $Z$ equals $c$;
for example, this occurs when all relative unitaries $U_iV_i^*$ agree.
\end{remark}

\section{Computational verification}

The reusable script \texttt{code/verify\_scaled\_unitary.py} generates seeded
complex Haar unitaries and compares direct Schatten norms with both formulas
\eqref{eq:full} and \eqref{eq:compressed}.  It also checks the inequality
direction for $p\in\{1,1.2,1.5,1.9,2,2.2,2.8,3.5,5,8\}$ across 5,000 random
instances with $1\leq d\leq6$ and $2\leq N\leq7$.  The exact command and
numerical residuals appear in \texttt{verification.md}.  These tests are not
used as proof.

\section{Limitations and bounded novelty check}

The theorem does not handle arbitrary blocks, blocks with non-flat singular
values, or the full matrix Hanner conjecture.  It is therefore promoted only
as a partial result.

The run's cheap indexes were searched by arXiv id, title, and the phrases
``norm compression'', ``Schatten'', and ``Hanner''.  A bounded web/arXiv
search combined the source title and author with ``unitary blocks'', ``scaled
unitaries'', and ``norm compression unitary''.  It recovered the source and
nearby PSD/diagonal results, but no exact statement of Theorem~\ref{thm:main}.
This is not an exhaustive novelty claim.

\section{Human-review recommendation}

The candidate is suitable for mathematical review.  The most useful checks
are: (i) the unitary equivalence following \eqref{eq:gram}; (ii) the
monotonicity sign in \eqref{eq:derivative}; and (iii) the factor $d$ in
\eqref{eq:compressed}.  Until review, retain the status
\emph{candidate partial result, likely valid}.

\section{Reference}

\begin{thebibliography}{9}
\bibitem{Audenaert}
K. M. R. Audenaert,
\emph{On a Norm Compression Inequality for $2\times N$ Partitioned Block
Matrices}, Linear Algebra Appl. 428 (2008), 781--795;
arXiv:math/0702186.
\end{thebibliography}

\end{document}
